\documentclass[../main.tex]{subfiles}
\begin{document}

\section{Matrix Decomposition}
Let $A \in \mathbb{R}^{m \times n}$ be the matrix that solves the system described by $Ax = b$ for $x \in \mathbb{R}^n$ and $b \in \mathbb{R}^m$.

\[
    A= \begin{bmatrix}
        a_{1,1} & a_{1,2} & \cdots & a_{1,n}\\
        a_{2,1} & a_{2,2} & \cdots & a_{2,n}\\
        \vdots  & \vdots  & \ddots & \vdots\\
        a_{m,1} & a_{m,2} & \cdots & a_{m,n}
    \end{bmatrix},
    \qquad
    x= \begin{bmatrix}
        x_1\\
        x_2\\
        \vdots\\
        x_n
    \end{bmatrix},
    \qquad
    b= \begin{bmatrix}
        b_1\\
        b_2\\
        \vdots\\
        b_m
    \end{bmatrix}.
\]

After performing \href{https://en.wikipedia.org/wiki/Gaussian_elimination}{Gaussian elimination} on $A$, $U$ (row reduction of $A$) is in row echelon form if:

\begin{itemize}
    \item All zero rows are at the bottom.
    \item The first nonzero entry in a row (from the left) is located to the right of the first nonzero entry in every row above.
\end{itemize}

The rank of $A$ is the number of nonzero rows in $U$. We say the system $A\mathbf{x} = \mathbf{b}$ is inconsistent (i.e. no solution) if $\text{rank}(A) < \text{rank}([A\mathbf{b}])$. The system has a unique solution when $\text{rank}(A) = \text{rank}([A\mathbf{b}]) = n$ (i.e. $\text{rank}(A)$ equals the number of variables in the system). The system has infinitely many solutions when $\text{rank}(A) = \text{rank}([A\mathbf{b}]) \text{ and rank}(A) < n$.

\subsection{LU Decomposition}

Any non-singular matrix $A$ can be decomposed into a product of a lower triangular matrix $L$ and an upper triangular matrix $U$ such that $A = LU$, using the procedures of Gaussian elimination.

\begin{thm}
    If $A$ can be reduced by Gaussian elimination to row echelon form using only operations of the form "add $c$ times row $j$ to row $i$" (that is, without scaling rows and without interchanging rows), then $A$ has an LU decomposition of the form $A = LU$, where $L$ is lower triangular and $U$ is upper triangular.

    In particular, after performing Gaussian elimination on $A$, the matrix $U$ is the corresponding row echelon form of $A$, and $L$ is given by
    \[
        L = \begin{pmatrix}
        1 & 0 & 0 & \cdots & 0 \\
        -c_{2,1} & 1 & 0 & \cdots & 0 \\
        -c_{3,1} & -c_{3,2} & 1 & \cdots & 0 \\
        \vdots & \vdots & \vdots & \ddots & \vdots \\
        -c_{m,1} & -c_{m,2} & -c_{m,3} & \cdots c& 1
        \end{pmatrix}
    \]
    where each entry corresponds to the elementary row operation "add $c_{i,j}$ times row $j$ to row $i$" performed during Gaussian elimination.
\end{thm}

If $A$ can be decomposed into $A = LU$, then $\mathbf{x}$ can be solved for by:

\begin{enumerate}
    \item Let $\mathbf{y} = U\mathbf{x}$ and solve the system for $L\mathbf{y} = \mathbf{b}$.
    \item Once $\mathbf{y}$ is found, the upper triangular system $U\mathbf{x} = \mathbf{y}$ can be solved with back substitution for $\mathbf{x}$.
\end{enumerate}

If we transform $A$ into $U$ by Gaussian elimination, the lower triangular matrix$L$ can be found by performing negations of the same operations on the identity matrix in reverse order. There are other methods to find $L$ and $U$, such as the Doolittle algorithm and the Crout algorithm.

\begin{table}[!htbp]
    \centering
    \renewcommand{\arraystretch}{1.4}
    \begin{tabular}{|c|p{0.82\textwidth}|}
        \hline
        \textbf{Steps} & \\ \hline
        1. & Initialize $L$ to an identity matrix, $\mathbb{I}$ of dimension $n\times n$, and $U=A$. \\ \hline
        2. & For $i=1,\ldots,n$ do Step 3. \\ \hline
        3. & \hspace{1.5em}For $j=i+1,\ldots,n$ do Steps 4--5. \\ \hline
        4. & \hspace{3em}Set $l_{ji}=u_{ji}/u_{ii}$. \\ \hline
        5. & \hspace{3em}Perform $U_j=(U_j-l_{ji}U_i)$ (where $U_i$, $U_j$ represent the $i$ and $j$ rows of the matrix $\mathbf{U}$, respectively). \\ \hline
    \end{tabular}
\end{table}

Once we have $L$ and $U$, we can solve for as many right-hand side vectors $\mathbf{b}$ as desired very quickly using a two-step process. First, let $\mathbf{y}=U\mathbf{x}$ and then solve $L\mathbf{y}=\mathbf{b}$ for $\mathbf{y}$ by forward substitution.
\subsubsection*{Pseudocode for forward substitution}
\begin{table}[!htbp]
    \centering
    \renewcommand{\arraystretch}{1.5}
    \begin{tabular}{|c|p{0.72\textwidth}|}
    \hline
    \textbf{Steps} & \\ \hline
    1. & Set $y_1=b_1/l_{11}$. \\ \hline
    2. & For $i=2,3,\ldots,n$ do Step 3. \\ \hline
    3. & \[
    y_i=\left(b_i-\sum_{j=1}^{i-1} l_{ij}y_j\right)\big/l_{ii}
    \] \\ \hline
    \end{tabular}
\end{table}
Then solve $U\mathbf{x}=\mathbf{y}$ for $\mathbf{x}$ using backward substitution.
\subsubsection*{Pseudocode for backward substitution}
\begin{table}[!htbp]
    \centering
    \renewcommand{\arraystretch}{1.5}
    \begin{tabular}{|c|p{0.72\textwidth}|}
    \hline
    \textbf{Steps} & \\ \hline
    1. & Set $x_n=y_n/u_{nn}$. \\ \hline
    2. & For $i=n-1,n-2,\ldots,1$ do Step 3. \\ \hline
    3. & \[
    x_i=\left(y_i-\sum_{j=i+1}^{n} u_{ij}x_j\right)\big/u_{ii}
    \] \\ \hline
    \end{tabular}
\end{table}

We can also easily solve for $A^{-1}$ and $\det(A)$ once $L$ and $U$ are known. For more details, see \href{https://johnfoster.pge.utexas.edu/numerical-methods-book/LinearAlgebra_LU.html}{Numerical Methods for Engineers}.


\noindent Sources:
\begin{itemize}
    \item \url{https://ubcmath.github.io/MATH307/systems/lu.html}
    \item \url{https://johnfoster.pge.utexas.edu/numerical-methods-book/LinearAlgebra_LU.html}
\end{itemize}

\end{document}